\documentclass[10pt,a4paper]{amsart}
\usepackage[margin=19mm]{geometry}
\usepackage{amsmath,amssymb,mathtools}
\usepackage[scr=rsfs,cal=euler]{mathalfa}
\usepackage{xfrac}
\usepackage[unicode,hidelinks,bookmarks=false]{hyperref}
\newcommand{\ie}{i.e.\ }
\newcommand{\cf}{cf.\ }
\newcommand{\resp}{respectively\ }
\newcommand{\Subsection}{\subsection}
\providecommand{\nobreakdash}{\nobreak\hbox{-}}
\setlength{\emergencystretch}{2em}
\multlinegap0pt
\usepackage{bm}
\usepackage{bbm}
\usepackage{dsfont}
\usepackage{xspace}
\usepackage{cancel}
\usepackage{mathtools}
\usepackage{amsthm}
\makeatletter
\@ifundefined{lemma}{\newtheorem{lemma}{Lemma}}{}
\makeatother
\newcommand{\Int}{\operatorname{int}}
\newcommand{\id}{\operatorname{id}}
\newcommand{\simeqp}{\simeq_{\mathrm{prop}}}
\title[Proper Homotopy Nonrigidity of Open Contractible Manifolds]{Proper Homotopy Nonrigidity of Open Contractible Manifolds}
\author{Donghan Kim}
\date{Source: arXiv:2607.18093v3 (CC BY 4.0)}
\begin{document}
\maketitle

\noindent\emph{Source and licence.} Proper Homotopy Nonrigidity of Open Contractible Manifolds. Version arXiv:2607.18093v3 (CC BY 4.0), distributed under the Creative Commons Attribution 4.0 International licence, as recorded in the arXiv licence metadata for this exact version and shown on the abstract page https://arxiv.org/abs/2607.18093v3. Full rights notice: licenses/arxiv-2607.18093-CC-BY-4.0.txt.\par\smallskip

\noindent\textit{Editorial scope.} Counted unit: the Lemma whose source label is \texttt{lem:interior-cone} in the article (source0.tex lines 441--447, source label \texttt{lem:interior-cone}), delivered together with the article's own complete original proof of it (source0.tex lines 449--512, one \texttt{proof} environment of 64 lines). No mathematics is written, repaired or re-derived: statement and proof are the article's own text line for line. Required context retained uncounted: none. Every definition, display and label of those lines is retained unchanged. Omitted from this package and not redistributed: 1--440, 448--448, 513--3606. Nothing inside a retained excerpt is omitted or altered: every retained body is delivered exactly as the article's lines are written, so no omission receipt of any kind accompanies this package. Citations: the retained excerpts carry no citation command at all, so no bibliography entry of the article is transcribed and no reference is redistributed. Reference adaptation: none was needed and none was made; every reference inside the delivered unit resolves inside it, so no unresolved reference or citation is produced. Printed slips kept literal: no printed word, symbol, hypothesis or number of the counted statement or of its proof was altered. Layout and class: the article's own class is \texttt{amsart} with packages including fontenc, inputenc, lmodern, microtype, amsmath, amssymb, amsthm, mathtools, mathrsfs, enumitem, xcolor, hyperref, tikz-cd, cleveref; the wrapper is the lane's pinned \texttt{amsart} editorial wrapper at 10pt on A4, which restates the theorem-like environments the retained text needs and adds the article's own macro definitions that the retained text needs (\texttt{\textbackslash Int}, \texttt{\textbackslash id}, \texttt{\textbackslash simeqp}), with extra wrapper packages (bm, bbm, dsfont, xspace, cancel, mathtools). The delivered numbering is the unit's own. Assets: no retained span contains a figure environment, a graphics-inclusion command, a PDF-page inclusion, a table or a TikZ picture; no image file is copied into this package, the package declares no asset dependency and no image file is redistributed with it. Licence: version arXiv:2607.18093v3 of this article is distributed under the Creative Commons Attribution 4.0 International licence, as recorded in the arXiv licence metadata for this exact version and shown on the abstract page https://arxiv.org/abs/2607.18093v3. A full rights notice is delivered at licenses/arxiv-2607.18093-CC-BY-4.0.txt. Characters: every retained excerpt is pure ASCII, so the delivered engine, XeLaTeX with its default Latin Modern text font, reports no missing character.
\begin{lemma}\label{lem:interior-cone}
Let \(C\) be a compact contractible manifold with nonempty connected
boundary \(Y=\partial C\).  Then
\[
\Int(C)\simeqp c^\circ Y.
\]
\end{lemma}

\begin{proof}
Choose a collar of the boundary and reparametrize its open part so
that
\[
\Int(C)=K\cup_Y\bigl(Y\times[0,\infty)\bigr),
\]
where \(K\) is a compact codimension-zero submanifold, the copy of
\(Y\) is identified with \(Y\times\{0\}\), and the inclusion
\(K\hookrightarrow C\) is a homotopy equivalence.  Hence \(K\) is
contractible.

Let
\[
q\colon \Int(C)\longrightarrow \Int(C)/K
\]
be the quotient map collapsing \(K\) to a point.  The quotient is
canonically homeomorphic to the open cone:
\[
\Int(C)/K\cong c^\circ Y.
\]
The inclusion \(K\hookrightarrow\Int(C)\) is a cofibration.  Since
\(K\) is contractible, it follows that \(q\) is a homotopy equivalence.

We now verify the proper statement.  Put
\[
T=Y\times[1,\infty)\subset\Int(C).
\]
Under the above identification of the quotient with \(c^\circ Y\),
the map \(q\) restricts to the identity on \(T\).  Both inclusions
\[
T\hookrightarrow\Int(C)
\qquad\text{and}\qquad
T\hookrightarrow c^\circ Y
\]
are cofibrations.  Since \(q\) is a homotopy equivalence and
\(q|_T=\id_T\), the relative homotopy-inverse theorem for
cofibrations, applied to
\[
q\colon(\Int(C),T)\longrightarrow(c^\circ Y,T),
\]
gives a homotopy inverse
\[
r\colon c^\circ Y\longrightarrow\Int(C)
\]
such that \(r|_T=\id_T\), together with homotopies
\[
rq\simeq\id_{\Int(C)}
\qquad\text{and}\qquad
qr\simeq\id_{c^\circ Y}
\]
which are stationary on \(T\).

The map \(q\) is proper because it collapses only the compact set
\(K\) and preserves the radial coordinate off \(K\).  The closures
of the complements of \(T\) in both spaces are compact.  Since \(r\)
is the identity on \(T\), the inverse image under \(r\) of a compact
set is a closed subset of the union of a compact radial core and a
compact subset of \(T\), and is therefore compact.  The inverse
homotopies are stationary on \(T\); the same compact-core argument
shows that they are proper.  Thus
\[
\Int(C)\simeqp c^\circ Y.
\]
\end{proof}

\end{document}
